% Generated draft: requires mathematical and source review.
\documentclass[11pt]{article}
\usepackage[margin=25mm]{geometry}
\usepackage{fontspec}
\setmainfont{Latin Modern Roman}
\usepackage{amsmath,amssymb,mathtools,unicode-math}
\setmathfont{Latin Modern Math}
\usepackage{enumitem,needspace,xurl,hyperref}
\setlength{\parindent}{0pt}
\setlength{\parskip}{5pt}
\setlength{\emergencystretch}{3em}
\newcommand{\N}{\mathbb{N}}
\newcommand{\Z}{\mathbb{Z}}
\newcommand{\Q}{\mathbb{Q}}
\newcommand{\R}{\mathbb{R}}
\newcommand{\C}{\mathbb{C}}
\newcommand{\F}{\mathbb{F}}
\newcommand{\A}{\mathbb{A}}
\newcommand{\PP}{\mathbb{P}}
\newcommand{\E}{\mathbb{E}}
\newcommand{\eps}{\varepsilon}
\newcommand{\dd}{\,\mathrm{d}}
\begin{document}
\section*{UnsolvedMath ID: 20002001}
\textbf{Source problem code:} catalogue entry
\subsection{Definitions and mathematical statement}
Use the minimum convention and the trivial valuation on \(\C\). For \(0\ne f=\sum_u c_ux^u\), let \(\operatorname{Trop}(f)\subset\R^n\) consist of weights where \(\min_{c_u\ne0}\langle w,u\rangle\) is attained by at least two distinct monomials. For an ideal \(I\), put \(\operatorname{Trop}(I)=\bigcap_{0\ne f\in I}\operatorname{Trop}(f)\). A finite subset \(J\subset I\) is a tropical basis if its hypersurface intersection equals this set; it need not be a universal Gr\"obner basis. The source asks to characterize or compute small such subsets, without fixing a size or computational-complexity objective.

It reports that the \(3\times3\) minors form a tropical basis for their determinantal ideal, and asks, through Sturmfels, whether the \(4\times4\) minors of a generic \(5\times5\) matrix do so. If \(I_4\subset\C[x_{ij}:1\leq i,j\leq5]\) is generated by those 25 minors, this asks
\[
\operatorname{Trop}(I_4)=\bigcap_{|R|=|C|=4}\operatorname{Trop}(\det X_{R,C}).
\]
The supplied aid reports an affirmative published answer by Chan--Jensen--Rubei, indeed for generic \(5\times N\) matrices. In rank language, all \(4\times4\) minors being tropically singular means tropical rank at most three, and the conclusion supplies a valued-field lift of classical rank at most three. Coefficient characteristic and valuation conventions must match the invoked correspondence theorem.

The aid's additional search reduction concerns the weight matrix \(A\in\R^{5\times5}\). Row and column potentials act by \(A_{ij}\mapsto A_{ij}+r_i+c_j\), with a one-dimensional redundancy in the potentials. Each orbit has the unique representative
\[
N(A)_{ij}=A_{ij}-A_{i1}-A_{1j}+A_{11},
\]
whose first row and first column vanish. This leaves 16 free entries. For every square minor, adding potentials shifts all its permutation costs by the same constant \(\sum_{i\in R}r_i+\sum_{j\in C}c_j\). The aid therefore reports preservation of the entire minimizing-permutation set and all cost gaps, including a zero first-to-second gap when the minimum is tied. In particular tropical singularity and tropical rank are unchanged.

Over an algebraically closed nonarchimedean field with elements realizing the required real valuations, multiplication of rows and columns of a lift by suitable nonzero scalars transports these potentials and preserves classical rank. Applying the inverse scaling transports lift certificates back. Thus Kapranov rank, defined as the smallest classical rank of a lift with those entry valuations, is preserved as well. The 16-coordinate normalization is a quotient of the certificate \emph{search space}; it is not a tropical basis consisting of 16 polynomials and does not by itself solve the smaller-basis question.
\subsection{Short English statement}
Characterize small collections of ideal polynomials whose tropical hypersurfaces define the whole tropical variety. The supplied aid reports the historical \(5\times5\) minor question answered affirmatively and gives a lossless 16-coordinate normalization for further searches. The normalization does not determine a smaller basis.
\subsection{Sources}
\begin{itemize}
\item Rekha Thomas, A.19, \emph{Tropical bases}, with the \(4\times4\) minor question attributed to Sturmfels, AIM workshop \emph{Amoebas and tropical geometry}: \url{https://aimath.org/WWN/amoebas/amoebas.pdf}.
\item The aid attributes the \(5\times N\) determinantal theorem to Chan--Jensen--Rubei without title, year or identifier; the original literature field is empty.
\item Local registry \texttt{data/unsolvedmath/problems.json}, permanent ID 20002001, code \texttt{catalogue entry}; canonical \texttt{aim-geometry-notes.json}, \texttt{notes[338]}; updated 14 August 2026.
\item Attempt 1: partial result, partially solved status, rigor label proved, novelty confidence low. The machine-generated aid is not peer reviewed and requires independent verification before citation.
\end{itemize}
\end{document}
